% 
% ###################################################################
% Copyright (c) 2018, AuthorName 
%  Editors: A.D. Rollett & M. De Graef
% All rights reserved.
%
% Licensed under the Creative Commons CC BY-NC-SA 4.0 License, 
% hereafter referred to as the "License"; you may not use this 
% document except in compliance with the License. You may obtain 
% a copy of the License at 
%     https://creativecommons.org/licenses/by-nc-sa/4.0/legalcode 
% Unless required by applicable law or agreed to in writing, all 
% material distributed under the License is distributed on an 
% "AS IS" BASIS, WITHOUT WARRANTIES OR CONDITIONS OF ANY KIND, 
% either express or implied. See the License for the specific 
% language governing permissions and limitations under the License.
% ###################################################################

% ###################################################################
% The following lines are to be uncommented or edited as needed 

%\corechapter{Yes}   % uncomment this line only if this chapter is a core/foundational chapter;
\OCchapterauthor{Marc De Graef, Carnegie Mellon University} %this will appear in a secondary header below the chapter title.

% graphics path for all figure *.eps files 

%	redefine the \chabbr command to be a six-character string that uniquely identifies this 
%	chapter; easiest way to do this is to take the first couple of letters from the title words
\renewcommand{\chabbr}{ORISYM}

%     replace \noheaderimage by the chapter header image file name (without .eps extension).
%     Chapter header images must be 2480 x 1240 pixels with 300dpi, RGB format.
\chapterimage{\noheaderimage}

%	Make sure to use the \label{\chabbr:} command everywhere instead of \label{} !
\chapter{Orientations and Symmetry}\label{\chabbr:OrSym}

% add the author information so that it will appear in the Author List at the start of the document
\writeauthor{ORISYM:OrSym}{Orientations and Symmetry}{De Graef}{Marc}{Materials Science and Engineering}{Carnegie Mellon University}{mdg@andrew.cmu.edu}{https://www.mse.engineering.cmu.edu/directory/bios/degraef-marc.html}

% Each chapter begins with Learning Objectives; the list of objectives should have links to sections/subsections using their OC labels
% each Learning Objective should be an active statement (i.e., contain a verb).  The \\ command can be used to force an item into the 
% second column if LaTeX breaks the line at an awkward location.
\begin{learningobjectives}
{\begin{itemize}
    \item[{\color{OCBurntOrange}\ref{\chabbr:sec:}:}] Learn 
    \item[{\color{OCBurntOrange}\ref{\chabbr:sec:}:}] Perform 
\end{itemize}}
\end{learningobjectives}

% ###################################################################
% ###################################################################
% ###################################################################

\section{Introductory Remarks}\label{\chabbr:sec:intro}
The orientation of an object with respect to an external reference frame can be described by the 3D rotation that is needed to bring this external reference frame into coincidence with the reference frame of the object.  













